\documentclass[11pt,a4paper]{article}
\usepackage[utf8]{inputenc}
\usepackage[T1]{fontenc}
\usepackage[catalan]{babel}
\usepackage{lmodern}
\usepackage[sfdefault,lining]{FiraSans}
\usepackage[a4paper,top=2.3cm,bottom=2.2cm,left=1.8cm,right=1.8cm,headheight=1.3cm,headsep=0.4cm,footskip=1cm]{geometry}
\usepackage{amsmath,amssymb}
\usepackage{icomma}
\usepackage{graphicx}
\usepackage[table]{xcolor}
\usepackage{tikz}
\usepackage{fontawesome5}
\usepackage{ccicons}
\usepackage{fancyhdr}
\usepackage{lastpage}
\usepackage{tasks}
\usepackage{needspace}
\usepackage[most]{tcolorbox}
\usepackage{array}
\usepackage[hidelinks,pdftitle={Fitxa d'exercicis de la unitat 3: Polinomis. 4t d'ESO Matemàtiques A},pdfsubject={Matemàtiques A 4t ESO},pdfkeywords={polinomis, monomis, valor numèric, suma de polinomis, resta de polinomis, producte de polinomis, identitats notables, divisió de polinomis, regla de Ruffini, teorema del residu, arrels, factorització, factor comú, fraccions algebraiques, simplificació}]{hyperref}

\definecolor{navy}{HTML}{1B2748}
\definecolor{accent}{HTML}{E8603A}
\definecolor{teal}{HTML}{1F9E90}
\definecolor{soft}{HTML}{EEF1F8}
\definecolor{muted}{HTML}{5D6685}

\setlength{\parindent}{0pt}
\setlength{\parskip}{3pt}
\settasks{label=\alph*), label-format=\bfseries\color{navy}, label-width=1.3em, item-indent=1.8em, column-sep=0.6em, after-item-skip=5pt, before-skip=3pt, after-skip=4pt}

% ---------- capçalera i peu ----------
\newcommand{\llicencia}{\href{https://creativecommons.org/licenses/by-nc-sa/4.0/deed.ca}{\ccbyncsa}\ \ Llicència Creative Commons Reconeixement-NoComercial-CompartirIgual 4.0 Internacional (CC BY-NC-SA 4.0)}
\pagestyle{fancy}
\fancyhf{}
\fancyhead[L]{\raisebox{-0.35\height}{\includegraphics[height=0.95cm]{logo-gva.png}}}
\fancyhead[C]{\small\color{muted}Fitxa d'exercicis · Unitat 3: Polinomis · 4t ESO A}
\fancyhead[R]{\raisebox{-0.35\height}{\includegraphics[height=0.95cm]{logo-vermellar.png}}}
\fancyfoot[L]{\scriptsize\color{muted}\llicencia}
\fancyfoot[R]{\small\color{muted}\thepage}
\renewcommand{\headrulewidth}{0.4pt}
\renewcommand{\headrule}{\color{muted!40}\hrule height \headrulewidth}
\fancypagestyle{primera}{\fancyhf{}\renewcommand{\headrulewidth}{0pt}
  \fancyfoot[L]{\scriptsize\color{muted}\llicencia}
  \fancyfoot[R]{\small\color{muted}\thepage}}

% ---------- nivells i exercicis ----------
\newcommand{\estrelles}[1]{{\color{accent}\ifcase#1\or\faStar\faStar[regular]\faStar[regular]\or\faStar\faStar\faStar[regular]\or\faStar\faStar\faStar\fi}}
\newcounter{exer}
\newcommand{\ex}[1]{\par\needspace{4\baselineskip}\medskip\refstepcounter{exer}\noindent{\color{navy}\bfseries\theexer.}\ \ {\small\estrelles{#1}}\ \ }

\newcommand{\seccio}[2]{\par\needspace{16\baselineskip}\bigskip
  \begin{tcolorbox}[colback=navy,colframe=navy,coltext=white,arc=2mm,boxrule=0pt,left=3mm,right=3mm,top=1.2mm,bottom=1.2mm]
  {\large\bfseries #1\quad #2}
  \end{tcolorbox}\vspace{-2pt}}

\newcommand{\exf}[1]{\par\medskip\refstepcounter{exer}\noindent{\color{navy}\bfseries\theexer.}\ \ {\small\estrelles{#1}}\ \ }
\newcolumntype{C}[1]{>{\centering\arraybackslash}p{#1}}
% quadre de recordatori (només a la fitxa d'exercicis)
\newcommand{\recorda}[1]{\par\smallskip\begin{tcolorbox}[colback=accent!7,colframe=accent!55,arc=2mm,boxrule=0.5pt,left=2.5mm,right=2.5mm,top=1mm,bottom=1mm]\small{\color{accent}\faLightbulb[regular]}\ \textbf{Recorda:}\ #1\end{tcolorbox}}
\newcommand{\res}[1]{#1}\newcommand{\rest}[1]{#1}

\def\fmneed{11}% que el títol d'una secció no quede sol al final del full
% =====================================================================
% Capçalera FentMates de les fitxes: un color per a cada curs.
% Els colors són els de la pàgina inicial (LEVELS a assets/js/materials.js);
% el logo ve de _eines/logo/pdf/ (python3 _eines/logo/genera.py).
%
% Ús, al final del preàmbul de la unitat:
%   % =====================================================================
% Capçalera FentMates de les fitxes: un color per a cada curs.
% Els colors són els de la pàgina inicial (LEVELS a assets/js/materials.js);
% el logo ve de _eines/logo/pdf/ (python3 _eines/logo/genera.py).
%
% Ús, al final del preàmbul de la unitat:
%   % =====================================================================
% Capçalera FentMates de les fitxes: un color per a cada curs.
% Els colors són els de la pàgina inicial (LEVELS a assets/js/materials.js);
% el logo ve de _eines/logo/pdf/ (python3 _eines/logo/genera.py).
%
% Ús, al final del preàmbul de la unitat:
%   \input{../comu/fentmates.tex}
%   \fmcurs{1eso}
% i, al començament del document:
%   \capcalera{Matemàtiques}{1r ESO}{Fitxa d'exercicis · Unitat 2}{Divisibilitat}
% Opcional abans de \input: \def\fmneed{14} (espai mínim abans de cada secció).
% =====================================================================
\makeatletter
\def\fm@logodir{../../logo/pdf}
\providecommand{\fmneed}{10}
\def\fm@cursnom{}

% clau del curs, color i color del text damunt del color
\newcommand{\fm@curs}[3]{\@namedef{fm@color@#1}{#2}\@namedef{fm@text@#1}{#3}}
\fm@curs{1eso}{E6B84A}{navy}
\fm@curs{2eso}{3B82F6}{white}
\fm@curs{3eso}{FF7A50}{white}
\fm@curs{4eso-a}{A78BFA}{white}
\fm@curs{4eso-b}{6D28D9}{white}
\fm@curs{1batx-mi}{9B1C1C}{white}
\fm@curs{1batx-ccss}{EF6B6B}{white}
\fm@curs{2batx-mii}{1E6B47}{white}
\fm@curs{2batx-ccss}{2BB3A3}{white}

\newcommand{\fmcurs}[1]{%
  \@ifundefined{fm@color@#1}{\PackageError{fentmates}{Curs desconegut: #1}{}}{}%
  \edef\fm@tmp{\noexpand\definecolor{curs}{HTML}{\@nameuse{fm@color@#1}}%
               \noexpand\colorlet{curstext}{\@nameuse{fm@text@#1}}}\fm@tmp
  \colorlet{cursfosc}{curs!72!black}%
  \def\fm@logo{\fm@logodir/vertical-#1.pdf}%
  \def\fm@marca{\fm@logodir/marca-#1.pdf}}

% ---------- capçalera de la primera pàgina ----------
% #1 matèria (petit), #2 curs, #3 tipus de document i unitat, #4 títol
\newcommand{\capcalera}[4]{%
  \gdef\fm@cursnom{#2}%
  \par\noindent
  \begin{tcolorbox}[colback=curs,colframe=curs,boxrule=0pt,arc=2.2mm,width=3.5cm,height=2.15cm,
      valign=center,left=3mm,right=2mm,top=1mm,bottom=1mm,nobeforeafter,box align=center]
    \color{curstext}{\footnotesize #1}\par\vspace{1pt}{\fontsize{19}{21}\selectfont\bfseries #2\par}%
  \end{tcolorbox}\hspace{5mm}%
  \begin{minipage}[c]{\dimexpr\linewidth-3.5cm-5mm-2.9cm\relax}
    {\color{cursfosc}\small\bfseries\MakeUppercase{#3}}\par\vspace{2pt}
    {\color{navy}\fontsize{20}{23}\selectfont\bfseries #4\par}
  \end{minipage}\hfill
  \begin{minipage}[c]{2.6cm}\raggedleft\includegraphics[width=2.6cm]{\fm@logo}\end{minipage}\par
  \vspace{2pt}}

% ---------- barres de secció del color del curs ----------
\renewcommand{\seccio}[2]{\par\needspace{\fmneed\baselineskip}\bigskip
  \begin{tcolorbox}[colback=curs,colframe=curs,coltext=curstext,arc=2mm,boxrule=0pt,left=3mm,right=3mm,top=1.2mm,bottom=1.2mm]
  {\large\bfseries #1\quad #2}
  \end{tcolorbox}\vspace{-2pt}}

% ---------- estrelles de dificultat del color del curs (a Batxillerat no se'n posen) ----------
\@ifundefined{estrelles}{}{%
  \renewcommand{\estrelles}[1]{{\color{curs}\ifcase#1\or\faStar\faStar[regular]\faStar[regular]\or\faStar\faStar\faStar[regular]\or\faStar\faStar\faStar\fi}}}

% ---------- capçalera i peu de les pàgines ----------
% primera pàgina: logos institucionals al peu
\fancypagestyle{primera}{\fancyhf{}\renewcommand{\headrulewidth}{0pt}
  \fancyfoot[L]{\raisebox{-0.32\height}{\includegraphics[height=0.78cm]{logo-gva.png}}\hspace{5mm}%
    \raisebox{-0.32\height}{\includegraphics[height=0.78cm]{logo-vermellar.png}}}
  \fancyfoot[R]{\scriptsize\color{muted}\href{https://creativecommons.org/licenses/by-nc-sa/4.0/deed.ca}{\ccbyncsa}\ \ CC BY-NC-SA 4.0 · fentmates.org\qquad\small\thepage}}
% la resta: el símbol a l'esquerra i el curs a la dreta, amb el títol al centre com abans
\fancyhead[L]{\raisebox{-0.3\height}{\includegraphics[height=0.62cm]{\fm@marca}}}
\fancyhead[R]{\small\bfseries\color{cursfosc}\fm@cursnom}
\renewcommand{\headrule}{\color{curs!45}\hrule height \headrulewidth}
\makeatother

%   \fmcurs{1eso}
% i, al començament del document:
%   \capcalera{Matemàtiques}{1r ESO}{Fitxa d'exercicis · Unitat 2}{Divisibilitat}
% Opcional abans de \input: \def\fmneed{14} (espai mínim abans de cada secció).
% =====================================================================
\makeatletter
\def\fm@logodir{../../logo/pdf}
\providecommand{\fmneed}{10}
\def\fm@cursnom{}

% clau del curs, color i color del text damunt del color
\newcommand{\fm@curs}[3]{\@namedef{fm@color@#1}{#2}\@namedef{fm@text@#1}{#3}}
\fm@curs{1eso}{E6B84A}{navy}
\fm@curs{2eso}{3B82F6}{white}
\fm@curs{3eso}{FF7A50}{white}
\fm@curs{4eso-a}{A78BFA}{white}
\fm@curs{4eso-b}{6D28D9}{white}
\fm@curs{1batx-mi}{9B1C1C}{white}
\fm@curs{1batx-ccss}{EF6B6B}{white}
\fm@curs{2batx-mii}{1E6B47}{white}
\fm@curs{2batx-ccss}{2BB3A3}{white}

\newcommand{\fmcurs}[1]{%
  \@ifundefined{fm@color@#1}{\PackageError{fentmates}{Curs desconegut: #1}{}}{}%
  \edef\fm@tmp{\noexpand\definecolor{curs}{HTML}{\@nameuse{fm@color@#1}}%
               \noexpand\colorlet{curstext}{\@nameuse{fm@text@#1}}}\fm@tmp
  \colorlet{cursfosc}{curs!72!black}%
  \def\fm@logo{\fm@logodir/vertical-#1.pdf}%
  \def\fm@marca{\fm@logodir/marca-#1.pdf}}

% ---------- capçalera de la primera pàgina ----------
% #1 matèria (petit), #2 curs, #3 tipus de document i unitat, #4 títol
\newcommand{\capcalera}[4]{%
  \gdef\fm@cursnom{#2}%
  \par\noindent
  \begin{tcolorbox}[colback=curs,colframe=curs,boxrule=0pt,arc=2.2mm,width=3.5cm,height=2.15cm,
      valign=center,left=3mm,right=2mm,top=1mm,bottom=1mm,nobeforeafter,box align=center]
    \color{curstext}{\footnotesize #1}\par\vspace{1pt}{\fontsize{19}{21}\selectfont\bfseries #2\par}%
  \end{tcolorbox}\hspace{5mm}%
  \begin{minipage}[c]{\dimexpr\linewidth-3.5cm-5mm-2.9cm\relax}
    {\color{cursfosc}\small\bfseries\MakeUppercase{#3}}\par\vspace{2pt}
    {\color{navy}\fontsize{20}{23}\selectfont\bfseries #4\par}
  \end{minipage}\hfill
  \begin{minipage}[c]{2.6cm}\raggedleft\includegraphics[width=2.6cm]{\fm@logo}\end{minipage}\par
  \vspace{2pt}}

% ---------- barres de secció del color del curs ----------
\renewcommand{\seccio}[2]{\par\needspace{\fmneed\baselineskip}\bigskip
  \begin{tcolorbox}[colback=curs,colframe=curs,coltext=curstext,arc=2mm,boxrule=0pt,left=3mm,right=3mm,top=1.2mm,bottom=1.2mm]
  {\large\bfseries #1\quad #2}
  \end{tcolorbox}\vspace{-2pt}}

% ---------- estrelles de dificultat del color del curs (a Batxillerat no se'n posen) ----------
\@ifundefined{estrelles}{}{%
  \renewcommand{\estrelles}[1]{{\color{curs}\ifcase#1\or\faStar\faStar[regular]\faStar[regular]\or\faStar\faStar\faStar[regular]\or\faStar\faStar\faStar\fi}}}

% ---------- capçalera i peu de les pàgines ----------
% primera pàgina: logos institucionals al peu
\fancypagestyle{primera}{\fancyhf{}\renewcommand{\headrulewidth}{0pt}
  \fancyfoot[L]{\raisebox{-0.32\height}{\includegraphics[height=0.78cm]{logo-gva.png}}\hspace{5mm}%
    \raisebox{-0.32\height}{\includegraphics[height=0.78cm]{logo-vermellar.png}}}
  \fancyfoot[R]{\scriptsize\color{muted}\href{https://creativecommons.org/licenses/by-nc-sa/4.0/deed.ca}{\ccbyncsa}\ \ CC BY-NC-SA 4.0 · fentmates.org\qquad\small\thepage}}
% la resta: el símbol a l'esquerra i el curs a la dreta, amb el títol al centre com abans
\fancyhead[L]{\raisebox{-0.3\height}{\includegraphics[height=0.62cm]{\fm@marca}}}
\fancyhead[R]{\small\bfseries\color{cursfosc}\fm@cursnom}
\renewcommand{\headrule}{\color{curs!45}\hrule height \headrulewidth}
\makeatother

%   \fmcurs{1eso}
% i, al començament del document:
%   \capcalera{Matemàtiques}{1r ESO}{Fitxa d'exercicis · Unitat 2}{Divisibilitat}
% Opcional abans de \input: \def\fmneed{14} (espai mínim abans de cada secció).
% =====================================================================
\makeatletter
\def\fm@logodir{../../logo/pdf}
\providecommand{\fmneed}{10}
\def\fm@cursnom{}

% clau del curs, color i color del text damunt del color
\newcommand{\fm@curs}[3]{\@namedef{fm@color@#1}{#2}\@namedef{fm@text@#1}{#3}}
\fm@curs{1eso}{E6B84A}{navy}
\fm@curs{2eso}{3B82F6}{white}
\fm@curs{3eso}{FF7A50}{white}
\fm@curs{4eso-a}{A78BFA}{white}
\fm@curs{4eso-b}{6D28D9}{white}
\fm@curs{1batx-mi}{9B1C1C}{white}
\fm@curs{1batx-ccss}{EF6B6B}{white}
\fm@curs{2batx-mii}{1E6B47}{white}
\fm@curs{2batx-ccss}{2BB3A3}{white}

\newcommand{\fmcurs}[1]{%
  \@ifundefined{fm@color@#1}{\PackageError{fentmates}{Curs desconegut: #1}{}}{}%
  \edef\fm@tmp{\noexpand\definecolor{curs}{HTML}{\@nameuse{fm@color@#1}}%
               \noexpand\colorlet{curstext}{\@nameuse{fm@text@#1}}}\fm@tmp
  \colorlet{cursfosc}{curs!72!black}%
  \def\fm@logo{\fm@logodir/vertical-#1.pdf}%
  \def\fm@marca{\fm@logodir/marca-#1.pdf}}

% ---------- capçalera de la primera pàgina ----------
% #1 matèria (petit), #2 curs, #3 tipus de document i unitat, #4 títol
\newcommand{\capcalera}[4]{%
  \gdef\fm@cursnom{#2}%
  \par\noindent
  \begin{tcolorbox}[colback=curs,colframe=curs,boxrule=0pt,arc=2.2mm,width=3.5cm,height=2.15cm,
      valign=center,left=3mm,right=2mm,top=1mm,bottom=1mm,nobeforeafter,box align=center]
    \color{curstext}{\footnotesize #1}\par\vspace{1pt}{\fontsize{19}{21}\selectfont\bfseries #2\par}%
  \end{tcolorbox}\hspace{5mm}%
  \begin{minipage}[c]{\dimexpr\linewidth-3.5cm-5mm-2.9cm\relax}
    {\color{cursfosc}\small\bfseries\MakeUppercase{#3}}\par\vspace{2pt}
    {\color{navy}\fontsize{20}{23}\selectfont\bfseries #4\par}
  \end{minipage}\hfill
  \begin{minipage}[c]{2.6cm}\raggedleft\includegraphics[width=2.6cm]{\fm@logo}\end{minipage}\par
  \vspace{2pt}}

% ---------- barres de secció del color del curs ----------
\renewcommand{\seccio}[2]{\par\needspace{\fmneed\baselineskip}\bigskip
  \begin{tcolorbox}[colback=curs,colframe=curs,coltext=curstext,arc=2mm,boxrule=0pt,left=3mm,right=3mm,top=1.2mm,bottom=1.2mm]
  {\large\bfseries #1\quad #2}
  \end{tcolorbox}\vspace{-2pt}}

% ---------- estrelles de dificultat del color del curs (a Batxillerat no se'n posen) ----------
\@ifundefined{estrelles}{}{%
  \renewcommand{\estrelles}[1]{{\color{curs}\ifcase#1\or\faStar\faStar[regular]\faStar[regular]\or\faStar\faStar\faStar[regular]\or\faStar\faStar\faStar\fi}}}

% ---------- capçalera i peu de les pàgines ----------
% primera pàgina: logos institucionals al peu
\fancypagestyle{primera}{\fancyhf{}\renewcommand{\headrulewidth}{0pt}
  \fancyfoot[L]{\raisebox{-0.32\height}{\includegraphics[height=0.78cm]{logo-gva.png}}\hspace{5mm}%
    \raisebox{-0.32\height}{\includegraphics[height=0.78cm]{logo-vermellar.png}}}
  \fancyfoot[R]{\scriptsize\color{muted}\href{https://creativecommons.org/licenses/by-nc-sa/4.0/deed.ca}{\ccbyncsa}\ \ CC BY-NC-SA 4.0 · fentmates.org\qquad\small\thepage}}
% la resta: el símbol a l'esquerra i el curs a la dreta, amb el títol al centre com abans
\fancyhead[L]{\raisebox{-0.3\height}{\includegraphics[height=0.62cm]{\fm@marca}}}
\fancyhead[R]{\small\bfseries\color{cursfosc}\fm@cursnom}
\renewcommand{\headrule}{\color{curs!45}\hrule height \headrulewidth}
\makeatother

\fmcurs{4eso-a}

\begin{document}
\thispagestyle{primera}

\capcalera{Matemàtiques A}{4t ESO}{Fitxa d'exercicis · Unitat 3}{\raggedright Polinomis}
\vspace{10pt}
Nom i cognoms: \rule{8.4cm}{0.4pt}\hfill Grup: \rule{1.6cm}{0.4pt}\hfill Data: \rule{2.2cm}{0.4pt}
\vspace{8pt}

\begin{tcolorbox}[colback=soft,colframe=soft,arc=2mm,left=3mm,right=3mm,top=1.5mm,bottom=1.5mm]
\small Exercicis graduats per nivells de dificultat:\quad \estrelles{1}\ bàsic \qquad \estrelles{2}\ mitjà \qquad \estrelles{3}\ avançat.\\[2pt] Resol-los al quadern mostrant tots els passos. Pots comprovar els resultats amb les eines interactives dels apunts.
\end{tcolorbox}

% =====================================================
\seccio{1}{Monomis}

\exf{1} Completa la taula:\begin{center}\renewcommand{\arraystretch}{1.5}\begin{tabular}{|C{2.4cm}|C{2.6cm}|C{2.6cm}|C{1.6cm}|}\hline\rowcolor{soft}\textbf{Monomi} & \textbf{Coeficient} & \textbf{Part literal} & \textbf{Grau}\\\hline$-3x^4$ &  &  & \\\hline$x^2y$ &  &  & \\\hline$\frac25x$ &  &  & \\\hline$-a^3b^2$ &  &  & \\\hline$8$ &  &  & \\\hline$-x$ &  &  & \\\hline\end{tabular}\end{center}

\ex{1} Agrupa els monomis semblants: \ $3x^2$, \ $-5x$, \ $x^3$, \ $7x^2$, \ $2x$, \ $-x^3$, \ $\frac12x^2$.

\ex{1} Redueix:\begin{tasks}(3)
  \task $5x+3x-x$
  \task $4x^2-7x^2+x^2$
  \task $2x^3+5x-x^3+3x$
  \task $3x^2-2x+5-x^2+4x-8$
  \task $\frac12x+\frac34x$
  \task $6xy-2xy+xy$
\end{tasks}\recorda{Només se sumen o resten els monomis semblants: se sumen els coeficients i es deixa la mateixa part literal.}

\ex{1} Calcula:\begin{tasks}(3)
  \task $3x^2\cdot4x^3$
  \task $(-2x)\cdot5x^4$
  \task $(-x^3)\cdot(-6x^2)$
  \task $\frac23x\cdot\frac94x^2$
  \task $4x^2y\cdot3xy^3$
  \task $2x\cdot(-3x)\cdot x^2$
\end{tasks}\recorda{En el producte, es multipliquen els coeficients i se sumen els exponents de cada lletra: $x^2\cdot x^3=x^5$.}

\ex{2} Calcula:\begin{tasks}(3)
  \task $12x^5:3x^2$
  \task $-20x^6:4x$
  \task $9x^4:(-3x^4)$
  \task $(2x^3)^2$
  \task $(-3x^2)^3$
  \task $\left(\frac12x\right)^4$
\end{tasks}

\ex{2} Calcula el valor numèric de cada monomi per als valors indicats:\begin{tasks}(1)
  \task $-2x^3$ \ per a \ $x=-1,\ 2,\ \frac{1}{2}$
  \task $\frac34x^2$ \ per a \ $x=-2,\ 0,\ 4$
\end{tasks}


% =====================================================
\seccio{2}{Polinomis i valor numèric}

\exf{1} Completa la taula:\begin{center}\renewcommand{\arraystretch}{1.5}\small\begin{tabular}{|l|C{3.7cm}|C{1.1cm}|C{1.5cm}|C{1.6cm}|C{1.95cm}|}\hline\rowcolor{soft}\textbf{Polinomi} & \textbf{Ordenat} & \textbf{Grau} & \textbf{Terme indep.} & \textbf{Coef. princ.} & \textbf{Complet?}\\\hline$P(x)=3x^2-5x^4+x-7$ &  &  &  &  & \\\hline$Q(x)=6-2x+x^3$ &  &  &  &  & \\\hline$R(x)=4x-x^2+1$ &  &  &  &  & \\\hline$S(x)=2x^5+x^2-3x$ &  &  &  &  & \\\hline$U(x)=\frac12x^3-x^2+3x-2$ &  &  &  &  & \\\hline\end{tabular}\end{center}

\ex{1} Escriu, en cada cas, un polinomi que complisca les condicions:\begin{tasks}(2)
  \task de grau 3, complet i amb terme independent $-2$;
  \task de grau 4 i amb només dos termes;
  \task de grau 2, amb coeficient principal $-3$;
  \task de grau 5 i sense terme independent.
\end{tasks}

\ex{1} Donat el polinomi $P(x)=x^3-2x^2+4x-1$, calcula $P(0)$, $P(1)$, $P(-1)$, $P(2)$ i $P(-2)$.\recorda{Substituïm la $x$ pel nombre, entre parèntesis si és negatiu: $(-2)^3=-8$, $(-2)^2=4$.}

\ex{1} Calcula el valor numèric dels polinomis següents per a $x=3$, $x=-1$ i $x=0$:\begin{tasks}(3)
  \task $P(x)=2x^{2}-5x+3$
  \task $P(x)=-x^{3}+4x$
  \task $P(x)=x^{4}-3x^{2}+2$
\end{tasks}

\ex{2} Quin valor ha de tindre $k$ perquè el polinomi $P(x)=2x^3-kx+5$ valga 9 per a $x=2$?

\ex{2} Troba $a$ en el polinomi $P(x)=x^2+ax-4$ si sabem que $P(-1)=-6$.

\ex{2} El cost, en euros, de fabricar $x$ samarretes és $C(x)=0,02x^2+3x+50$. Quant costa fabricar-ne 10? I 100?

\ex{3} L'altura (en metres) d'una pilota llançada cap amunt és $h(t)=-5t^2+20t+1$, on $t$ és el temps en segons. Calcula $h(0)$, $h(1)$, $h(2)$ i $h(4)$. Què signifiquen $h(0)$ i $h(4)$?


% =====================================================
\seccio{3}{Suma i resta de polinomis}

\exf{1} Donats $P(x)=3x^3-2x^2+x-5$, $Q(x)=x^3+4x^2-6$ i $R(x)=-2x^2+3x+1$, calcula:\begin{tasks}(2)
  \task $P(x)+Q(x)$
  \task $P(x)-Q(x)$
  \task $Q(x)-R(x)$
  \task $P(x)+Q(x)+R(x)$
\end{tasks}\recorda{Per a restar, canviem el signe de \emph{tots} els termes del polinomi que restem: $-(x^2-3x)=-x^2+3x$.}

\ex{2} Amb els mateixos polinomis, calcula:\begin{tasks}(3)
  \task $P(x)-[Q(x)+R(x)]$
  \task $2P(x)-3R(x)$
  \task $P(x)-2Q(x)+R(x)$
\end{tasks}

\ex{2} Troba el polinomi $Q(x)$ tal que $\left(x^3-2x+4\right)+Q(x)=3x^3+x^2-5$.

\ex{2} Els costats d'un triangle fan $x+3$, $2x-1$ i $x^2+x$ centímetres. Escriu el polinomi que dona el perímetre i calcula'l per a $x=4$.


% =====================================================
\seccio{4}{Producte de polinomis}

\exf{1} Multiplica:\begin{tasks}(2)
  \task $3x\left(x^2-2x+5\right)$
  \task $-2x^2\left(x^3-4x+1\right)$
  \task $\frac12x\left(4x^2-6x+2\right)$
  \task $5x^3\left(-x^2+x-3\right)$
  \task $\left(x^2-3x+1\right)\cdot(-4x)$
  \task $-x\left(2x^4-x^2+7\right)$
\end{tasks}\recorda{Es multiplica el monomi per cada terme del polinomi. Compte amb els signes: $(-)\cdot(-)=+$.}

\ex{1} Multiplica i redueix:\begin{tasks}(2)
  \task $(x+2)(x-5)$
  \task $(2x-1)(3x+4)$
  \task $\left(x^2+1\right)(x-3)$
  \task $\left(x^2-2x+3\right)(x+1)$
  \task $\left(2x^2-x+4\right)(x-2)$
  \task $\left(x^2+x-1\right)\left(x^2-x+2\right)$
\end{tasks}\recorda{Cada terme del primer polinomi es multiplica per cada terme del segon; després es redueixen els termes semblants.}

\ex{2} Donats $P(x)=x^2-3x+2$, $Q(x)=2x+1$ i $R(x)=x-4$, calcula:\begin{tasks}(3)
  \task $P(x)\cdot Q(x)$
  \task $P(x)\cdot R(x)-Q(x)$
  \task $[Q(x)+R(x)]\cdot P(x)$
\end{tasks}

\ex{2} Opera i redueix:\begin{tasks}(1)
  \task $2x(x-3)-x(x+1)$
  \task $(x+2)(x-1)-(x-3)(x+4)$
  \task $3\left(x^2-1\right)-(x+1)(2x-3)$
\end{tasks}

\ex{3} Un rectangle fa $2x+3$ cm de base i $x-1$ cm d'altura. Escriu el polinomi que dona l'àrea i calcula-la per a $x=5$. I el volum d'una caixa d'arestes $x$, $x+2$ i $x+5$? Calcula'l per a $x=3$.


% =====================================================
\seccio{5}{Identitats notables}

\exf{1} Desenvolupa aplicant les identitats notables:\begin{tasks}(3)
  \task $(x+3)^2$
  \task $(x-4)^2$
  \task $(x+5)(x-5)$
  \task $(2x+1)^2$
  \task $(3x-2)^2$
  \task $(2x+7)(2x-7)$
  \task $\left(x^2+3\right)^2$
  \task $(5-x)^2$
  \task $\left(\frac12x+1\right)^2$
\end{tasks}\recorda{$(a+b)^2=a^2+2ab+b^2$, \quad $(a-b)^2=a^2-2ab+b^2$, \quad $(a+b)(a-b)=a^2-b^2$.}

\ex{2} Opera i simplifica:\begin{tasks}(2)
  \task $(x+1)^2-(x-1)^2$
  \task $(x+3)^2-(x+3)(x-3)$
  \task $(2x-1)^2+(x+2)(x-2)$
  \task $2(x-3)^2-(x+1)^2$
\end{tasks}

\ex{2} Completa els buits:\begin{tasks}(1)
  \task $(x+\square)^2=x^2+\square+36$
  \task $(\square-5)^2=4x^2-\square+25$
  \task $(3x+\square)(3x-\square)=\square-16$
\end{tasks}

\ex{3} Calcula sense calculadora, fent servir les identitats notables: \ $101^2$, \ $99^2$ \ i \ $51\cdot49$.


% =====================================================
\seccio{6}{Divisió de polinomis}

\exf{1} Divideix:\begin{tasks}(1)
  \task $\left(6x^3-9x^2+3x\right):3x$
  \task $\left(10x^5-15x^3+5x^2\right):5x^2$
  \task $\left(8x^4-4x^3+2x^2\right):\left(-2x^2\right)$
\end{tasks}

\ex{2} Fes les divisions següents i indica el quocient i el residu:\begin{tasks}(1)
  \task $\left(2x^3+3x^2-x+4\right):\left(x^2+1\right)$
  \task $\left(x^4-2x^3+3x-1\right):\left(x^2-x\right)$
  \task $\left(6x^3-x^2+x+5\right):(2x-1)$
  \task $\left(x^4-1\right):\left(x^2+1\right)$
  \task $\left(3x^4+2x^2-x\right):\left(x^2-2x+1\right)$
\end{tasks}\recorda{Ordena el dividend i deixa buit el lloc dels termes que falten. Para quan el residu tinga grau menor que el divisor.}

\ex{2} Comprova el resultat de la divisió a) de l'exercici anterior amb la prova de la divisió, $D(x)=d(x)\cdot c(x)+r(x)$.

\ex{3} En una divisió, el divisor és $x^2-2$, el quocient és $x+3$ i el residu és $2x-1$. Quin és el dividend?

\ex{3} L'àrea d'un rectangle és $2x^2+7x+3$ i la base fa $x+3$. Quant fa l'altura? Quin és el perímetre?


% =====================================================
\seccio{7}{Regla de Ruffini}

\exf{1} Divideix aplicant la regla de Ruffini i indica el quocient i el residu:\begin{tasks}(2)
  \task $\left(x^{3}-4x^{2}+x+6\right):(x-2)$
  \task $\left(x^{3}+2x^{2}-x+3\right):(x+1)$
  \task $\left(2x^{3}-3x^{2}+4\right):(x-1)$
  \task $\left(x^{4}-3x^{2}+2x-5\right):(x+2)$
  \task $\left(x^{4}-16\right):(x-2)$
  \task $\left(3x^{3}-5x+1\right):(x+3)$
\end{tasks}\recorda{Escriu tots els coeficients, amb un 0 per a cada terme que falte. Si el divisor és $x+2$, a l'esquerra va el $-2$.}

\ex{2} Divideix per Ruffini:\begin{tasks}(2)
  \task $\left(x^5-1\right):(x-1)$
  \task $\left(x^3+27\right):(x+3)$
\end{tasks}

\ex{2} Fes servir la regla de Ruffini per a saber si $x^3-2x^2-5x+6$ és divisible entre $x-3$ i entre $x+1$.

\ex{3} Troba el valor de $k$ perquè la divisió $\left(x^3+kx^2-4x+12\right):(x-2)$ siga exacta.


% =====================================================
\seccio{8}{Valor numèric amb Ruffini. Teorema del residu}

\exf{1} Calcula, sense fer la divisió, el residu de:\begin{tasks}(2)
  \task $\left(x^{3}-5x+2\right):(x-2)$
  \task $\left(x^{4}+2x^{3}-x\right):(x+1)$
  \task $\left(2x^{5}-3x^{2}+1\right):(x-1)$
  \task $\left(x^{3}+x^{2}+x+1\right):(x+2)$
\end{tasks}\recorda{El residu de la divisió $P(x):(x-a)$ és el valor numèric $P(a)$.}

\ex{2} Donat $P(x)=x^4-2x^3-5x^2+x+3$, calcula $P(3)$ i $P(-2)$ amb la regla de Ruffini. Comprova-ho substituint.

\ex{2} Troba $m$ perquè el residu de la divisió $\left(x^3-mx+4\right):(x+2)$ siga 10.

\ex{3} Quin valor ha de tindre $a$ perquè el polinomi $2x^3+ax^2-x+3$ done residu 5 en dividir-lo entre $x-1$?


% =====================================================
\seccio{9}{Arrels d'un polinomi}

\exf{1} Comprova si $x=1$, $x=-1$ i $x=2$ són arrels del polinomi $P(x)=x^3-2x^2-x+2$.

\ex{1} Quins d'aquests nombres són arrels de $x^3+x^2-6x$? \ $0$, \ $2$, \ $-2$, \ $-3$.

\ex{2} Troba les arrels enteres dels polinomis:\begin{tasks}(2)
  \task $x^{3}-6x^{2}+11x-6$
  \task $x^{3}+2x^{2}-x-2$
  \task $x^{3}-3x+2$
  \task $x^{4}-5x^{2}+4$
\end{tasks}\recorda{Les arrels enteres són divisors del terme independent: prova'ls amb Ruffini (o calculant el valor numèric).}

\ex{2} Troba les arrels dels polinomis de segon grau resolent l'equació:\begin{tasks}(2)
  \task $x^{2}-7x+12$
  \task $2x^{2}-x-1$
  \task $x^{2}+2x+5$
  \task $x^{2}-6x+9$
\end{tasks}

\ex{2} Sense fer cap càlcul, explica per què $x=5$ no pot ser una arrel de $x^3+x^2-3x+6$.

\ex{3} Escriu un polinomi de grau 3, amb coeficient principal 1, que tinga per arrels $2$, $-1$ i $3$. Troba també $k$ perquè $x=2$ siga arrel de $x^3+kx-6$.


% =====================================================
\seccio{10}{Factor comú i identitats notables}

\exf{1} Trau factor comú:\begin{tasks}(3)
  \task $4x^{2}-8x$
  \task $3x^{3}+6x^{2}-9x$
  \task $5x^{4}-10x^{3}$
  \task $12x^{5}-18x^{3}+6x^{2}$
  \task $x^{3}+x^{2}+x$
  \task $2x^{4}-14x^{2}$
\end{tasks}\recorda{Es trau el m.c.d. dels coeficients i la menor potència de $x$: $6x^3-9x^2=3x^2(2x-3)$.}

\ex{1} Factoritza fent servir les identitats notables:\begin{tasks}(3)
  \task $x^{2}-36$
  \task $x^{2}+8x+16$
  \task $x^{2}-12x+36$
  \task $9x^{2}-4$
  \task $4x^{2}+20x+25$
  \task $25x^{2}-10x+1$
  \task $x^{4}-9$
  \task $x^{2}-\frac{1}{4}$
  \task $16-x^2$
\end{tasks}\recorda{$a^2+2ab+b^2=(a+b)^2$, \quad $a^2-2ab+b^2=(a-b)^2$, \quad $a^2-b^2=(a+b)(a-b)$.}

\ex{2} Factoritza: primer trau factor comú i després aplica una identitat notable.\begin{tasks}(3)
  \task $2x^{3}-18x$
  \task $3x^{3}+12x^{2}+12x$
  \task $x^{4}-x^{2}$
  \task $5x^{3}-10x^{2}+5x$
  \task $x^{5}-16x$
  \task $2x^{4}-50x^{2}$
\end{tasks}

\ex{3} Calcula sense calculadora, factoritzant: \ $83^2-17^2$.


% =====================================================
\seccio{11}{Factorització de polinomis}

\exf{2} Factoritza aquests polinomis de segon grau (busca'n primer les arrels):\begin{tasks}(2)
  \task $x^{2}-5x+6$
  \task $x^{2}+x-12$
  \task $2x^{2}-8x+6$
  \task $x^{2}-4x-5$
\end{tasks}\recorda{Si $x_1$ i $x_2$ són les arrels de $ax^2+bx+c$, aleshores $ax^2+bx+c=a(x-x_1)(x-x_2)$.}

\ex{2} Factoritza amb la regla de Ruffini:\begin{tasks}(2)
  \task $x^{3}-2x^{2}-5x+6$
  \task $x^{3}+4x^{2}+x-6$
  \task $x^{3}-7x-6$
  \task $x^{3}-3x^{2}-4x+12$
\end{tasks}

\ex{3} Factoritza (si es pot, trau primer factor comú):\begin{tasks}(2)
  \task $x^{4}-x^{3}-7x^{2}+x+6$
  \task $2x^{3}-4x^{2}-2x+4$
  \task $x^{4}-2x^{3}-3x^{2}$
  \task $x^{3}-x^{2}-8x+12$
\end{tasks}

\ex{3} Factoritza i indica les arrels i quantes vegades es repeteix cadascuna:\begin{tasks}(2)
  \task $x^{3}-3x^{2}+3x-1$
  \task $x^{4}-8x^{2}+16$
\end{tasks}


% =====================================================
\seccio{12}{Aplicació: simplificació de fraccions algebraiques}

\exf{1} Simplifica traient factor comú en el numerador i en el denominador:\begin{tasks}[after-item-skip=8pt](3)
  \task $\dfrac{6x^{3}}{9x^{2}}$
  \task $\dfrac{4x^{2}-8x}{4x}$
  \task $\dfrac{3x+6}{x^{2}+2x}$
  \task $\dfrac{x^{2}-5x}{2x-10}$
  \task $\dfrac{5x^{3}+10x^{2}}{15x^{2}}$
  \task $\dfrac{2x-6}{4x-12}$
\end{tasks}\recorda{Per a simplificar una fracció algebraica, \textbf{factoritzem} el numerador i el denominador i eliminem els \textbf{factors} comuns: $\dfrac{x^2+3x}{5x}=\dfrac{x(x+3)}{5x}=\dfrac{x+3}{5}$.}

\ex{2} Simplifica fent servir les identitats notables (i el factor comú, si cal):\begin{tasks}[after-item-skip=8pt](3)
  \task $\dfrac{x^{2}-9}{x+3}$
  \task $\dfrac{x^{2}-4}{x^{2}+4x+4}$
  \task $\dfrac{2x^{2}-2}{x^{2}-2x+1}$
  \task $\dfrac{x^{2}+6x+9}{2x+6}$
  \task $\dfrac{x^{3}-x}{x^{2}+x}$
  \task $\dfrac{3x^{2}-12}{3x^{2}-6x}$
\end{tasks}

\ex{3} Simplifica (factoritza amb l'equació de segon grau o amb Ruffini):\begin{tasks}[after-item-skip=8pt](2)
  \task $\dfrac{x^{2}-5x+6}{x^{2}-9}$
  \task $\dfrac{x^{2}+x-2}{x^{2}-1}$
  \task $\dfrac{x^{3}-2x^{2}-5x+6}{x^{2}+x-2}$
  \task $\dfrac{x^{3}-4x}{x^{3}-4x^{2}+4x}$
\end{tasks}

\ex{2} Hi ha alumnes que han simplificat així. Indica quines simplificacions són correctes i explica l'error de les altres.\begin{tasks}[after-item-skip=8pt](4)
  \task $\dfrac{x+6}{6}=x$
  \task $\dfrac{3x}{6x^2}=\dfrac{1}{2x}$
  \task $\dfrac{x^2-1}{x-1}=x+1$
  \task $\dfrac{x^2+4}{x+2}=x+2$
\end{tasks}


% =====================================================
\seccio{13}{Problemes}

\exf{2} A un quadrat de costat $x+2$ li retallem un quadrat de costat $x$. Escriu el polinomi que dona l'àrea que queda i calcula-la per a $x=3$ cm.

\ex{2} D'una làmina de cartó de $20\times30$ cm retallem un quadrat de costat $x$ en cada cantó i pleguem les vores per a fer una caixa sense tapa. Escriu el polinomi $V(x)$ del volum de la caixa i calcula'l per a $x=5$ cm.

\ex{3} L'àrea d'un rectangle és $x^2+5x+6$ i un dels costats fa $x+2$. Quant fa l'altre costat? Escriu el perímetre i calcula'l per a $x=4$.

\ex{3} Pensa un nombre, suma-li 3, eleva el resultat al quadrat i resta-li el quadrat del nombre pensat. Comprova amb polinomis que el resultat sempre és el séxtuple del nombre més 9.

\vfill
\begin{center}\color{muted}\small\faCheckCircle[regular]\ \ Ànim! Comprova els teus resultats amb les eines interactives dels apunts de la unitat.\end{center}

\end{document}
